\Article{3}
Wenn $m$ eine complexe ganze Zahl bedeutet, die mit $\mu$ keinen gemeinschaftlichen Divisor hat, und die sämmtlichen Zahlen eines Viertelsystems $\kappa$ mit $m$ multiplicirt werden, so bilden die Producte, oder beliebige ihnen congruente Zahlen ihrerseits auch ein Viertelsystem; und ebenso entstehen durch Multiplication der Zahlen der Systeme $\kappa',\kappa'',\kappa'''$ noch drei Viertelsysteme, die mit jenem zusammengehören werden. Der Beweis ist leicht zu führen. Diese vier neuen Systeme mögen mit $m\kappa,m\kappa',m\kappa'',m\kappa'''$ bezeichnet werden, gleichviel, ob die Producte selbst oder nur ihnen congruente Zahlen gewählt werden. Im letztern Fall kann dies so geschehen, dass man immer nur solche wählt, die sich in einem der Systeme $\kappa,\kappa',\kappa'',\kappa'''$ befinden. Auf diese Art ist also das System $\chi$, wenigstens allgemein zu reden, auf zwei verschiedene Arten in Viertelsysteme zerlegt. Nehmen wir an, dass $m\kappa$ gemeinschaftlich hat

\[
\begin{array}{rl}
\text{mit }\kappa\dotsfive&\lambda\text{ Zahlen}\\
\kappa'\dotsfive&\lambda'\text{ Zahlen}\\
\kappa''\dotsfive&\lambda''\text{ Zahlen}\\
\kappa'''\dotsfive&\lambda'''\text{ Zahlen}
\end{array}
\]
so wird auch $\kappa'$ mit $m\kappa'$, $\kappa''$ mit $m\kappa''$, $\kappa'''$ mit $m\kappa'''$ gemein haben $\lambda$ Glieder;
\SourceLeaf{380}
$\kappa''$ mit $m\kappa'$, $\kappa'''$ mit $m\kappa''$, $\kappa$ mit $m\kappa'''$, $\lambda'$ Glieder u.\,s.\,w. Es sei $\varepsilon$ der kleinste Rest von $\lambda'+2\lambda''+3\lambda'''$ nach dem Modulus $4$, oder $\varepsilon$ eine der vier Zahlen $0,1,2,3$, je nachdem $\lambda'+2\lambda''+3\lambda'''$ von der Form $4n,4n+1,4n+2,4n+3$ ist. Ich behaupte nun, dass $\varepsilon$ von der Anordnung des Viertelsystems $\kappa$ unabhängig ist.

Um die Beweisführung zu erleichtern, bediene ich mich folgender Bezeichnung. $\Pi\psi$ soll die Zahl $0,1,2,3$ bezeichnen, je nachdem die durch $\mu$ nicht theilbare Zahl $\psi$ sich (selbst oder durch Congruenz Repräsentation) in der Gruppe $\kappa,\kappa',\kappa'',\kappa'''$ befindet. Von selbst hat man daher die Folge

I.\quad $\Pi(i\psi)\equiv1+\Pi\psi\ (\text{mod. }4)$.

II.\quad Die Zahl $i^{-\Pi\psi}\pdot\psi$ findet sich, entweder selbst oder durch Congruenz Repräsentation in der Gruppe $\kappa$.

III.\quad $\Sigma\Pi m\varphi\equiv\varepsilon,\ (\text{mod. }4)$, wenn die Summation über alle in $\kappa$ befindliche Glieder $\varphi$ erstreckt wird.

Es sei nun $k$ ein anderes Viertelsystem, bestehend aus $f,f',f''$ u.\,s.\,w., während $\kappa$ aus $\varphi,\varphi',\varphi''$ u.\,s.\,w. besteht. Ich setze voraus, was erlaubt ist, da die \emph{Ordnung} der Glieder in $\kappa$ willkürlich, dass $f$ mit $\varphi$ identisch oder zusammenhängend ist, $f'$ mit $\varphi'$, $f''$ mit $\varphi''$ u.\,s.\,w. Die mit $k$ zusammenhängenden Viertelsgruppen seien $k'(\equiv ik)$, $k''(\equiv-k)$, $k'''(\equiv-ik)$. Es habe ferner die Charakteristik $\mathrm{P}$ in Beziehung auf die Gruppen $k,k',k'',k'''$ dieselbe Bedeutung wie $\Pi$ in Beziehung auf $\kappa,\kappa',\kappa'',\kappa'''$, so dass $\mathrm{P}\psi=0,1,2,3$, je nachdem $\psi$ zu $k,k',k'',k'''$ gehört.

Es wird demnach, wenn man von der Vertheilung der $\chi$ in die Viertelsysteme $k,k',k'',k'''$ anstatt von $\kappa,\kappa',\kappa'',\kappa'''$ ausgeht, an die Stelle von $\varepsilon$ treten der kleinste Rest von $\mathrm{P}mf+\mathrm{P}mf'+\mathrm{P}mf''+\mathrm{P}mf'''$ u.\,s.\,w. u.\,s.\,w. oder von $\Sigma\mathrm{P}mf$ nach dem Modulus $4$ und es handelt sich, zu beweisen, dass $\Sigma\mathrm{P}mf-\Sigma\Pi m\varphi$ durch $4$ theilbar ist.

Wir schreiben diese Grösse so
\[
\begin{array}{l}
\phantom{+}\mathrm{P}mf+\mathrm{P}mf'+\mathrm{P}mf''+\mathrm{P}mf'''+\text{u.\,s.\,w.}\\
-\mathrm{P}m\varphi-\mathrm{P}m\varphi'-\mathrm{P}m\varphi''-\mathrm{P}m\varphi'''-\\
+\mathrm{P}m\varphi+\mathrm{P}m\varphi'+\mathrm{P}m\varphi''+\mathrm{P}m\varphi'''+\\
-\Pi m\varphi-\Pi m\varphi'-\Pi m\varphi''-\Pi m\varphi'''-
\end{array}
\]

Da der Voraussetzung nach $f$ und $\varphi$ congruent sind oder zusammengehören, so gilt dasselbe auch von $mf$ und $m\psi$ und man hat

\SourceLeaf{381}
\[
\left.\begin{aligned}
f&\equiv i^{-\mathrm{P}\varphi}\varphi\\
i^{-\mathrm{P}mf}mf&\equiv i^{-\mathrm{P}m\varphi}m\varphi
\end{aligned}\right\}\ \text{mod. }\mu
\]
woraus leicht folgt $\mathrm{P}mf-\mathrm{P}m\varphi\equiv-\mathrm{P}\varphi\ (\text{mod. }4)$ und das Aggregat der beiden obersten Reihen $\equiv-\Sigma\mathrm{P}\varphi$. Da nun ferner $\mathrm{P}m\varphi-\Pi m\varphi\equiv\mathrm{P}(i^{-\Pi m\varphi}m\varphi)$ ist, $m\varphi\pdot i^{-\Pi m\varphi}$ zu $\kappa$ gehört und der Inbegriff \emph{aller} $m\varphi\pdot i^{-\Pi m\varphi}$ ohne Rücksicht auf die Ordnung mit allen $\varphi$ übereinkommt, so wird das Aggregat \emph{aller} $\mathrm{P}(m\varphi\pdot i^{-\Pi m\varphi})$ aequal sein dem Aggregat aller $\mathrm{P}\varphi$; folglich das Aggregat der dritten und vierten Reihe $\equiv\Sigma\mathrm{P}\varphi\ (\text{mod. }4)$, also das Aggregat aller vier Reihen $\equiv0\ (\text{mod. }4)$ W.\ Z.\ B.\ W.

Da also $\varepsilon$, unabhängig von der Wahl der Viertelsysteme blos von $m$ und $\mu$ abhängt, so werden wir $\varepsilon$ den Character der Zahl $m$ in Beziehung auf den Modulus $\mu$ nennen und mit $\text{Ch. }m\ (\text{mod. }\mu)$ bezeichnen. Man sieht leicht, dass dies nur eine Generalisirung derjenigen Definition ist, die (Art.\ $\ldots$) für den Fall, wo $\mu$ eine Primzahl ist, gegeben ist, oder sie unter sich begreift.
